\documentclass{article}
\usepackage[UTF8]{ctex}
\usepackage{amsmath,amssymb,amsfonts}
\usepackage{enumitem}
\usepackage{geometry}
\geometry{a4paper, margin=1in}
% === 错题本排版优化（不参与题目标识） ===
\setlength{\parindent}{0pt}
\setlength{\parskip}{0.35em plus 0.10em minus 0.08em}
\setlength{\emergencystretch}{3em}
\setlist{itemsep=0.25em, topsep=0.40em, parsep=0pt}
\allowdisplaybreaks[2]
% === 错题本排版优化结束 ===
\begin{document}
\section*{东北师范大学2024年高等代数真题}
\section*{71．东北师范大学2024年高等代数与解析几何试卷}

\subsection{题目1}
1．（15 分）求下列行列式

\[
\left|\begin{array}{cccc}
1+x_{1} & 1+x_{1}^{2} & \cdots & 1+x_{1}^{n} \\
1+x_{2} & 1+x_{2}^{2} & \cdots & 1+x_{2}^{n} \\
\vdots & & & \\
1+x_{n} & 1+x_{n}^{2} & \cdots & 1+x_{n}^{n}
\end{array}\right|
\]

\subsection{题目2}
2．（15 分）设 \(A\) 为一个 4 阶实对称矩阵，特征值为0，0，0，4，属于特征值0的线性无关的特征向量为

\[
(-1,1,0,0)^{\prime},(-1,0,1,0)^{\prime},(-1,0,0,1)^{\prime} .
\]

求矩阵 \(A\) 。

\subsection{题目3}
3．（10 分）设 \(f(x)=x^{m}-1, g(x)=x^{n}-1\) ，其中 \(m, n\) 为正整数，证明：\((f(x), g(x))=x^{d}-1\) ，其中 \(d\) 为 \(m, n\) 的最大公因式。

\subsection{题目4}
4．（15 分）设 \(V_{1}, V_{2}, \cdots, V_{m}\) 是有理数域上向量空间 \(V\) 的 \(m\) 个真子空间，证明：\(V\) 中必定有一个向量\(\alpha, \alpha\) 不属于任何一个 \(V_{i}\) 。

\subsection{题目5}
5．（15 分）设 \(A\) 是复数域上的 \(n\) 阶幂等矩阵，且 \(A=A_{1}+A_{2}+\cdots+A_{k}\) ，其中 \(k\) 是大于等于 2 的正整数，又有 \(r(A)=r\left(A_{1}\right)+r\left(A_{2}\right)+\cdots+r\left(A_{k}\right)\) ．证明：\(A_{i}^{2}=A_{i}, A_{i} A_{j}=O(i \neq j)\) ．

\subsection{题目6}
6．（15 分）如果二次型 \(f\left(x_{1}, x_{2}, \cdots, x_{n}\right)\) 仅在 \(x_{1}=x_{2}=\cdots=x_{n}=0\) 时为零，证明：\(f\) 必定是正定型或负定型。

\subsection{题目7}
7．（15 分）设 \(\varphi\) 是 \(n\) 维欧氏空间 \(V\) 上的对称变换，证明：对 \(V\) 中任意向量 \(\alpha\) 都有 \((\varphi(\alpha), \alpha) \geq 0\) 的充分必要条件是 \(\varphi\) 的全部特征值都为非负实数。

\subsection{题目8}
8．（20 分）设直线 \(l_{1}: \frac{x-1}{2}=\frac{y}{1}=\frac{z}{-1}, l_{2}: \frac{x}{2}=\frac{y}{1}=\frac{z+1}{-2}\) 和平面 \(\pi: x+y+z=0\) ，求与直线 \(l_{1}, l_{2}\) 相交且平行于 \(\pi\) 的所有直线构成的曲面 \(S\) 的方程．

\subsection{题目9}
9．（20 分）证明方程 \(5 x^{2}+5 y^{2}+2 z^{2}-8 x y-2 x z-2 y z+20 x+20 y-40 z-16=0\) 表示的曲面是柱面，并指出此柱面的母线方向及一条准线方程。

\subsection{题目10}
10．（10 分）证明：\(O, A, B\) 不共线，\(A \neq B\) ，则点 \(C\) 与点 \(A, B\) 共线的充分必要条件是向量 \(\overrightarrow{O C}\) 可以分解为向量 \(\overrightarrow{O A}\) 和向量 \(\overrightarrow{O B}\) 的线性组合，且分解的系数之和为1。
\end{document}
